\documentclass[12pt,a4paper]{ctexart}
\usepackage[top=2.5cm,bottom=2.5cm,left=2.5cm,right=2.5cm]{geometry}
\usepackage{amsmath,amssymb,amsfonts}
\usepackage{graphicx}
\usepackage{tikz}
\usetikzlibrary{shapes,arrows,positioning,calc,decorations.markings,matrix,patterns,fit,backgrounds}
\usepackage{xcolor}
\usepackage{enumitem}
\usepackage{booktabs}
\usepackage{hyperref}
\usepackage{physics}
\usepackage{float}
\usepackage{caption}
\usepackage{subcaption}
\usepackage{framed}
\usepackage{tcolorbox}
\tcbuselibrary{skins,breakable}
\usepackage{algorithm}
\usepackage{algorithmic}

\hypersetup{
    colorlinks=true,
    linkcolor=blue!70!black,
    citecolor=green!50!black,
    urlcolor=blue!70!black
}

% 定义颜色
\definecolor{tensorblue}{RGB}{70,130,180}
\definecolor{tensorred}{RGB}{205,92,92}
\definecolor{tensorgreen}{RGB}{60,179,113}
\definecolor{tensororange}{RGB}{230,160,50}
\definecolor{tensorpurple}{RGB}{147,112,219}
\definecolor{queencolor}{RGB}{220,50,50}
\definecolor{swapcolor}{RGB}{255,165,0}
\definecolor{ctmcolor}{RGB}{100,149,237}
\definecolor{envcolor}{RGB}{144,238,144}
\definecolor{darkgreen}{RGB}{0,100,0}

% TikZ 样式
\tikzset{
    tensor/.style={draw, thick, minimum size=0.6cm},
    sitetensor/.style={draw, thick, fill=tensorblue!20, minimum size=0.7cm},
    swaptensor/.style={draw, thick, fill=swapcolor!30, minimum size=0.5cm, diamond},
    cornertensor/.style={draw, thick, fill=ctmcolor!30, minimum size=0.6cm},
    edgetensor/.style={draw, thick, fill=envcolor!30, minimum size=0.6cm},
    cleg/.style={thick, tensorblue},
    pleg/.style={thick, tensorred},
    hbond/.style={thick, tensorblue},
    vbond/.style={thick, tensorgreen},
    dbond/.style={thick, tensorred},
    abond/.style={thick, tensororange},
    queen/.style={fill=queencolor!80, text=white, font=\bfseries\large},
    empty/.style={fill=gray!10},
}

\floatname{algorithm}{算法}
\renewcommand{\algorithmicrequire}{\textbf{输入:}}
\renewcommand{\algorithmicensure}{\textbf{输出:}}

\title{\textbf{$N$ 皇后问题的张量网络深度分析\\——从非平面拓扑到二维收缩方法}}
\author{}
\date{}

\begin{document}
\maketitle
\thispagestyle{empty}

\begin{abstract}
本文是对 $N$ 皇后问题张量网络方法的\textbf{深度分析}，超越了基础的 MPS-MPO 逐行收缩框架。我们从张量网络的\textbf{拓扑结构}出发，指出 $N$ 皇后网络因对角线约束产生的键交叉而本质上是\textbf{非平面的}——这一关键特征使得标准的二维张量网络方法（TRG、CTMRG 等）不能直接适用。我们提出通过引入 SWAP 张量进行\textbf{平面化}的策略，将非平面网络转化为可用二维方法处理的平面网络。在此基础上，我们详细讨论了 TRG（张量重正化群）、HOTRG（高阶 TRG）和 CTMRG（角转移矩阵重正化群）三种二维收缩方法在 $N$ 皇后问题上的适用性与实现细节。我们进一步分析了 $N$ 皇后问题的\textbf{纠缠结构}，论证了精确计数时键维度指数增长的必然性（\#P-完全性的体现），并探讨了皇后格点气体在热力学极限下的\textbf{相变行为}。最后，我们给出了四种不同计算方案的具体实现方案，并与蒙特卡洛方法进行了系统对比。
\end{abstract}

\tableofcontents
\newpage

%=============================================================
\section{引言与回顾}
%=============================================================

\subsection{为什么需要深度分析？}

在前文《$N$ 皇后问题的张量网络解法——从约束编码到精确计数》中，我们建立了一套完整的方法链：将 $N$ 皇后问题的约束（行、列、对角线）编码为矩阵乘积态（MPS），融合得到格点张量，构造转移矩阵并分解为矩阵乘积算符（MPO），最终通过 MPS-MPO 逐行收缩计算解的总数 $Z_N$。

这套方法在概念上是清晰完整的，但它本质上是一个\textbf{一维方法应用于二维问题}。当我们把棋盘按行展开、逐行处理时，二维的空间关联被压缩到了一维的 MPS 表示中。这种压缩的代价是 MPS 键维度 $\chi$ 的指数增长——对于 $N = 8$，精确键维度在中间行达到 568；对于更大的 $N$，$\chi$ 按 $O(2^N/\sqrt{N})$ 增长。

一个自然的问题是：\textbf{能否保持二维结构，用真正的二维张量网络方法来处理这个问题？}

答案并不简单。$N$ 皇后问题的张量网络具有一个微妙但关键的特征：由于对角线约束产生的键交叉，它是一个\textbf{非平面网络}。这意味着标准的二维张量网络方法（如用于 Ising 模型的 TRG 或 CTMRG）不能直接适用——它们都假设网络可以嵌入到平面上。

本文的目标是深入分析这一问题的拓扑结构，提出解决方案，并给出可操作的计算方案。

\subsection{1D MPS-MPO 方法回顾}

为了使本文自包含，我们简要回顾 MPS-MPO 方法的核心框架。

在 $N \times N$ 棋盘上，为每个格点 $(i,j)$ 引入二值变量 $n_{ij} \in \{0, 1\}$（1 表示放置皇后）。$N$ 皇后解的总数可以表示为配分函数：
\begin{equation}
Z_N = \sum_{\{n_{ij}\}} \prod_{i=1}^{N} \delta\!\left(\sum_j n_{ij} = 1\right) \cdot \prod_{j=1}^{N} \delta\!\left(\sum_i n_{ij} = 1\right) \cdot \prod_{d \in \mathcal{D}} \delta\!\left(\sum_{(i,j)\in d} n_{ij} \leq 1\right)
\end{equation}
其中三个乘积分别对应行约束（每行恰好一个）、列约束（每列恰好一个）和对角线约束（每条对角线至多一个）。

\textbf{逐行处理的思路}：利用每行恰好一个皇后的约束，用变量 $q_i \in \{1, \ldots, N\}$ 表示第 $i$ 行皇后所在的列。定义增广状态 $|\mathbf{s}\rangle = |\mathbf{c}, \mathbf{r}, \mathbf{l}\rangle$，其中：
\begin{itemize}
\item $\mathbf{c} = (c_1, \ldots, c_N) \in \{0,1\}^N$：列占据向量（$c_j = 1$ 表示列 $j$ 已被占据）
\item $\mathbf{r} = (r_1, \ldots, r_N) \in \{0,1\}^N$：右对角线（``$\backslash$''方向）威胁向量
\item $\mathbf{l} = (l_1, \ldots, l_N) \in \{0,1\}^N$：左对角线（``/''方向）威胁向量
\end{itemize}

状态空间维度为 $2^{3N}$，但转移矩阵 $T$ 可以分解为 MPO，键维度仅为 $D_{\text{MPO}} = 4$（与 $N$ 无关）。配分函数为：
\begin{equation}
\boxed{Z_N = \langle \text{final} | \, T^N \, | \text{initial} \rangle}
\end{equation}
通过 MPS-MPO 逐行乘法计算。精确 MPS 键维度在 $m \approx N/2$ 时达到峰值 $\chi_{\max} \sim \binom{N}{N/2}$。

\subsection{本文的结构}

本文按以下逻辑展开：

\begin{itemize}
\item \textbf{第 2 节}：深入分析 1D 方法的局限性——非平面拓扑与指数键维度增长
\item \textbf{第 3 节}：8-腿格点张量的详细解剖——四类键的几何与交叉问题
\item \textbf{第 4 节}：平面化策略——SWAP 张量的引入与融合
\item \textbf{第 5 节}：2D 收缩方法——TRG、HOTRG 与 CTMRG
\item \textbf{第 6 节}：皇后格点气体的热力学极限与相变
\item \textbf{第 7 节}：纠缠结构与计算复杂度的深层联系
\item \textbf{第 8 节}：与蒙特卡洛方法的对比
\item \textbf{第 9 节}：对称性的利用——$D_4$ 群
\item \textbf{第 10 节}：四种具体计算方案
\item \textbf{第 11 节}：讨论与展望
\end{itemize}


%=============================================================
\section{一维方法的深层局限}
%=============================================================

\subsection{非平面性：一个被忽视的拓扑障碍}

让我们从几何角度重新审视 $N$ 皇后的张量网络。每个格点 $(i,j)$ 上的张量 $\mathcal{T}$ 有 8 条腿，对应 4 种约束链：

\begin{center}
\begin{tikzpicture}[scale=1.2]
  % 中心节点
  \node[sitetensor, minimum size=1cm] (T) at (0,0) {$\mathcal{T}$};

  % 水平键（行约束）
  \draw[hbond, ->] (-2,0) -- (T.west) node[midway, above] {$\alpha_L$};
  \draw[hbond, ->] (T.east) -- (2,0) node[midway, above] {$\alpha_R$};

  % 垂直键（列约束）
  \draw[vbond, ->] (0,2) -- (T.north) node[midway, right] {$\beta_U$};
  \draw[vbond, ->] (T.south) -- (0,-2) node[midway, right] {$\beta_D$};

  % "/"对角线键
  \draw[abond, ->] (-1.5,-1.5) -- (T.south west) node[midway, left] {$\gamma_{\text{in}}$};
  \draw[abond, ->] (T.north east) -- (1.5,1.5) node[midway, right] {$\gamma_{\text{out}}$};

  % "\"对角线键
  \draw[dbond, ->] (-1.5,1.5) -- (T.north west) node[midway, left] {$\delta_{\text{in}}$};
  \draw[dbond, ->] (T.south east) -- (1.5,-1.5) node[midway, right] {$\delta_{\text{out}}$};

  % 图例
  \node[right] at (3, 1.5) {\textcolor{tensorblue}{--- 行约束（水平）}};
  \node[right] at (3, 0.8) {\textcolor{tensorgreen}{--- 列约束（垂直）}};
  \node[right] at (3, 0.1) {\textcolor{tensororange}{--- ``/''对角线}};
  \node[right] at (3, -0.6) {\textcolor{tensorred}{--- ``$\backslash$''对角线}};
\end{tikzpicture}
\end{center}

当我们把所有格点张量按棋盘排列并连接相应的键时，水平键和垂直键构成一个标准的方格子——这是平面的，没有任何问题。但对角线键的加入引入了\textbf{交叉}。

\begin{tcolorbox}[title=关键观察：对角线键的交叉, colback=red!5, colframe=red!60!black]
考虑棋盘上的一个基本方块（plaquette），由四个格点 $(i,j)$、$(i,j+1)$、$(i+1,j)$、$(i+1,j+1)$ 构成。在这个方块中：
\begin{itemize}
\item ``$\backslash$''方向：格点 $(i,j)$ 的 $\delta_{\text{out}}$ 连接到 $(i+1,j+1)$ 的 $\delta_{\text{in}}$
\item ``/''方向：格点 $(i,j+1)$ 的 $\gamma_{\text{out}}$ 连接到 ... 不对，``/''方向是 $i+j = \text{const}$，即 $(i,j+1)$ 连接到 $(i+1,j)$... 但这些在同一行，不会直接相连。
\end{itemize}

更准确地说：``$\backslash$''对角线从 $(i,j)$ 到 $(i+1,j+1)$（左上到右下），``/''对角线从 $(i,j)$ 到 $(i+1,j-1)$（右上到左下）。在每个方块内部，从 $(i,j)$ 到 $(i+1,j+1)$ 的 ``$\backslash$''键与从 $(i,j+1)$ 到 $(i+1,j)$ 的 ``/''键\textbf{必然交叉}。
\end{tcolorbox}

让我们用一个 $3 \times 3$ 的例子来看清这一点：

\begin{center}
\begin{tikzpicture}[scale=1.8]
  % 格点
  \foreach \i in {0,1,2} {
    \foreach \j in {0,1,2} {
      \node[sitetensor, minimum size=0.5cm] (\i-\j) at (\j, -\i) {};
    }
  }

  % 水平键（蓝色）
  \foreach \i in {0,1,2} {
    \draw[hbond] (\i-0) -- (\i-1);
    \draw[hbond] (\i-1) -- (\i-2);
  }

  % 垂直键（绿色）
  \foreach \j in {0,1,2} {
    \draw[vbond] (0-\j) -- (1-\j);
    \draw[vbond] (1-\j) -- (2-\j);
  }

  % "\"对角线键（红色）
  \draw[dbond] (0-0) -- (1-1);
  \draw[dbond] (0-1) -- (1-2);
  \draw[dbond] (1-0) -- (2-1);
  \draw[dbond] (1-1) -- (2-2);

  % "/"对角线键（橙色）
  \draw[abond] (0-1) -- (1-0);
  \draw[abond] (0-2) -- (1-1);
  \draw[abond] (1-1) -- (2-0);
  \draw[abond] (1-2) -- (2-1);

  % 标记交叉点
  \fill[red!50, opacity=0.5] (0.5, -0.5) circle (0.12);
  \fill[red!50, opacity=0.5] (1.5, -0.5) circle (0.12);
  \fill[red!50, opacity=0.5] (0.5, -1.5) circle (0.12);
  \fill[red!50, opacity=0.5] (1.5, -1.5) circle (0.12);

  % 标注
  \node[right, red!70!black] at (2.5, -1) {红色圆圈 = 键交叉点};
  \node[right] at (2.5, -0.3) {\textcolor{tensorblue}{蓝} = 行，\textcolor{tensorgreen}{绿} = 列};
  \node[right] at (2.5, -1.7) {\textcolor{tensorred}{红} = $\backslash$对角线，\textcolor{tensororange}{橙} = /对角线};
\end{tikzpicture}
\end{center}

在 $3 \times 3$ 的网络中，有 $(3-1)^2 = 4$ 个键交叉点。一般地，$N \times N$ 网络有 $(N-1)^2$ 个键交叉。

\textbf{为什么非平面性是一个问题？}几乎所有成熟的二维张量网络收缩算法——TRG、HOTRG、CTMRG、SRG 等——都假设网络是\textbf{平面的}。平面性保证了我们可以定义一致的``内部''和``外部''，可以从边界开始逐步收缩网络。非平面网络没有这样的全局结构，标准方法无法直接使用。

\subsection{一维压缩的代价：纠缠的视角}

MPS-MPO 方法将二维网络``展平''为一维。这意味着原本在二维中是短程的关联，在一维表示中可能变成长程的。

具体来说，考虑 MPS 在处理完前 $m$ 行后的状态 $|\psi_m\rangle$。这个 MPS 的键维度 $\chi_m$ 等于``在给定前 $m$ 行所有可能的合法部分解的前提下，第 $m$ 行和第 $m+1$ 行之间需要传递的信息量''。

这个信息量包括：
\begin{enumerate}
\item \textbf{列占据信息}：前 $m$ 行中有 $m$ 列已被占据，具体是哪 $m$ 列？有 $\binom{N}{m}$ 种可能。
\item \textbf{右对角线威胁}：从前 $m$ 行的皇后延伸下来的 ``$\backslash$'' 对角线攻击。
\item \textbf{左对角线威胁}：从前 $m$ 行的皇后延伸下来的 ``/'' 对角线攻击。
\end{enumerate}

仅列占据信息一项就给出了键维度的下界：
\begin{equation}
\boxed{\chi_m \geq \binom{N}{m}}
\end{equation}

这在 $m = N/2$ 时给出 $\chi_{N/2} \geq \binom{N}{N/2} \sim \frac{2^N}{\sqrt{\pi N/2}}$。加上对角线信息后，实际的键维度更大。

\begin{center}
\begin{tabular}{c|cccccccccc}
\toprule
$N$ & 4 & 6 & 8 & 10 & 12 & 14 & 16 & 18 & 20 \\
\midrule
$\binom{N}{N/2}$ & 6 & 20 & 70 & 252 & 924 & 3432 & 12870 & 48620 & 184756 \\
$\chi_{N/2}$（精确） & 6 & 36 & 568 & -- & -- & -- & -- & -- & -- \\
\bottomrule
\end{tabular}
\end{center}

实际的键维度（考虑对角线信息后）比 $\binom{N}{N/2}$ 还要大数倍。

\subsection{面积律 vs. 体积律}

在凝聚态物理中，纠缠熵有一个重要的分类标准：

\begin{itemize}
\item \textbf{面积律（area law）}：二分纠缠熵 $S$ 正比于分割面的``面积''（在一维中是常数，在二维中正比于边长）。物理基态通常满足面积律，MPS/PEPS 可以高效表示。
\item \textbf{体积律（volume law）}：$S$ 正比于子系统的``体积''。随机态、高温态满足体积律，需要指数大的键维度。
\end{itemize}

$N$ 皇后问题的解空间具有一种\textbf{类似体积律}的纠缠结构。原因在于\textbf{列约束的全局性}：``所有列恰好使用一次''本质上等价于一个\textbf{排列}约束。已知排列群 $S_N$ 的表示论给出，在中间切割处的 Schmidt 秩（即不同的列占据模式数）为 $\binom{N}{N/2}$——这是体积律级别的增长。

\begin{tcolorbox}[title=与 Ising 模型的对比, colback=blue!5, colframe=tensorblue]
\textbf{二维 Ising 模型}：
\begin{itemize}
\item 格点张量：4 条腿（水平 + 垂直），键维度 $D = 2$
\item 网络拓扑：标准方格子，\textbf{平面}
\item 纠缠结构：面积律（临界点处有对数修正）
\item MPS-MPO 键维度：$\chi \sim O(1)$（远离临界点）或 $O(\text{poly}(N))$（临界点附近）
\item CTMRG 环境维度 $\chi \sim 10\text{--}100$ 已足够精确
\end{itemize}

\textbf{$N$ 皇后问题}：
\begin{itemize}
\item 格点张量：8 条腿（水平 + 垂直 + 两个对角线），键维度 $D = 2$
\item 网络拓扑：方格子 + 对角线，\textbf{非平面}
\item 纠缠结构：体积律级别（$S \sim N$）
\item MPS-MPO 键维度：$\chi \sim O(2^N/\sqrt{N})$（指数增长）
\item 精确收缩需要指数资源
\end{itemize}
\end{tcolorbox}

\subsection{为什么二维方法仍然有价值？}

既然精确计数注定指数困难（\#P-完全性的必然结果），为什么还要发展二维方法？

\begin{enumerate}
\item \textbf{更好的近似效率}：同样的截断键维度 $\chi$，二维方法可以捕获更多的空间关联，因为它不需要把二维信息压缩到一维。
\item \textbf{热力学极限}：对于皇后格点气体，CTMRG 可以直接在 $N \to \infty$ 下计算每格点自由能——这是一维方法完全无法做到的。
\item \textbf{相变研究}：CTMRG 自然地给出关联长度、序参量等物理量，可以精确定位相变点。
\item \textbf{物理直觉}：二维方法保持了问题的空间结构，使得物理图像更加直观。
\end{enumerate}


%=============================================================
\section{8-腿格点张量的深入解剖}
%=============================================================

\subsection{格点张量的完整定义}

让我们完整地写出格点张量的表达式。在格点 $(i,j)$ 上，四条约束链共享同一个物理变量 $n_{ij} \in \{0, 1\}$。格点张量通过对 $n_{ij}$ 求和得到：
\begin{equation}
\mathcal{T}_{\alpha_L \alpha_R,\, \beta_U \beta_D,\, \gamma_{\text{in}} \gamma_{\text{out}},\, \delta_{\text{in}} \delta_{\text{out}}} = \sum_{n=0}^{1} A^{[n]}_{\alpha_L \alpha_R} \cdot A^{[n]}_{\beta_U \beta_D} \cdot B^{[n]}_{\gamma_{\text{in}} \gamma_{\text{out}}} \cdot B^{[n]}_{\delta_{\text{in}} \delta_{\text{out}}}
\end{equation}
其中 $A^{[n]}$ 是``恰好一个''约束的局部张量，$B^{[n]}$ 是``至多一个''约束的局部张量：
\begin{equation}
A^{[0]} = \begin{pmatrix} 1 & 0 \\ 0 & 1 \end{pmatrix} = I, \quad
A^{[1]} = \begin{pmatrix} 0 & 1 \\ 0 & 0 \end{pmatrix}
\end{equation}
\begin{equation}
B^{[0]} = \begin{pmatrix} 1 & 0 \\ 0 & 1 \end{pmatrix} = I, \quad
B^{[1]} = \begin{pmatrix} 0 & 1 \\ 0 & 0 \end{pmatrix}
\end{equation}

注意 $A^{[n]}$ 和 $B^{[n]}$ 的矩阵形式完全相同！区别在于\textbf{边界条件}：行/列约束链的右端（下端）向量为 $|R\rangle = (0,1)^T$（要求恰好一个），对角线约束链的右端向量为 $|R'\rangle = (1,1)^T$（允许零个或一个）。

\subsection{非零元素的枚举}

由于 $n$ 只取 0 和 1，格点张量只有两个"贡献项"：

\textbf{$n = 0$（不放皇后）}：
\begin{equation}
\mathcal{T}^{(n=0)} = I \otimes I \otimes I \otimes I = \delta_{\alpha_L \alpha_R}\, \delta_{\beta_U \beta_D}\, \delta_{\gamma_{\text{in}} \gamma_{\text{out}}}\, \delta_{\delta_{\text{in}} \delta_{\text{out}}}
\end{equation}
这是 $2^4 = 16$ 个非零元素（4 个指标对各自独立，每对有 2 种取值），即所有 $(\alpha, \alpha, \beta, \beta, \gamma, \gamma, \delta, \delta)$ 形式的元素值为 1。

\textbf{物理含义}：不放皇后时，四条约束链的内部状态不变地``穿过''这个格点。

\textbf{$n = 1$（放皇后）}：
\begin{equation}
\mathcal{T}^{(n=1)} = A^{[1]} \otimes A^{[1]} \otimes B^{[1]} \otimes B^{[1]}
\end{equation}
由于 $A^{[1]}$ 和 $B^{[1]}$ 只有一个非零元素 $(A^{[1]})_{01} = 1$，所以 $\mathcal{T}^{(n=1)}$ 也只有一个非零元素：
\begin{equation}
\mathcal{T}_{01,\, 01,\, 01,\, 01} = 1
\end{equation}
即所有四个指标对都必须从 0 跳到 1。

\textbf{物理含义}：放皇后时，四条约束链必须\textbf{同时}从``未见皇后''（状态 0）跳到``已见皇后''（状态 1）。如果任何一条链已经处于状态 1（该行/列/对角线已有皇后），则不可能再跳到状态 1——这就自动地实施了互斥约束。

\textbf{总计}：格点张量有 $16 + 1 = 17$ 个非零元素（在 $2^8 = 256$ 个元素中），非零比例仅为 $17/256 \approx 6.6\%$。这种\textbf{极端稀疏性}是设计高效算法时应当利用的重要特征。

\subsection{四类键的几何分析}

让我们仔细分析四类键在 $N \times N$ 棋盘上的连接方式：

\textbf{水平键（行约束，蓝色）}：格点 $(i,j)$ 的 $\alpha_R$ 连接到 $(i,j+1)$ 的 $\alpha_L$。方向：从左到右。每行 $i$ 有 $N-1$ 条水平键，加上两端的边界向量 $\langle L| = (1,0)$（左端）和 $|R\rangle = (0,1)^T$（右端），共同实施``每行恰好一个皇后''的约束。

\textbf{垂直键（列约束，绿色）}：格点 $(i,j)$ 的 $\beta_D$ 连接到 $(i+1,j)$ 的 $\beta_U$。方向：从上到下。每列 $j$ 有 $N-1$ 条垂直键。

\textbf{``$\backslash$''对角线键（主对角线方向，红色）}：格点 $(i,j)$ 的 $\delta_{\text{out}}$ 连接到 $(i+1,j+1)$ 的 $\delta_{\text{in}}$。方向：从左上到右下。同一条 ``$\backslash$'' 对角线上 $i - j = \text{const}$ 的所有格点被串联。

\textbf{``/''对角线键（反对角线方向，橙色）}：格点 $(i,j)$ 的 $\gamma_{\text{out}}$ 连接到 $(i+1,j-1)$ 的 $\gamma_{\text{in}}$。方向：从右上到左下。同一条 ``/'' 对角线上 $i + j = \text{const}$ 的所有格点被串联。

\subsection{交叉的精确计数}

在棋盘的每个基本方块（plaquette）中，存在恰好一个键交叉——``$\backslash$'' 键从左上角到右下角，``/'' 键从右上角到左下角，这两条键在方块中心交叉。

\begin{center}
\begin{tikzpicture}[scale=2.5]
  % 四个格点
  \node[sitetensor, minimum size=0.6cm] (TL) at (0,1) {};
  \node[sitetensor, minimum size=0.6cm] (TR) at (1,1) {};
  \node[sitetensor, minimum size=0.6cm] (BL) at (0,0) {};
  \node[sitetensor, minimum size=0.6cm] (BR) at (1,0) {};

  % 标签
  \node[above] at (TL.north) {\small $(i,j)$};
  \node[above] at (TR.north) {\small $(i,j\!+\!1)$};
  \node[below] at (BL.south) {\small $(i\!+\!1,j)$};
  \node[below] at (BR.south) {\small $(i\!+\!1,j\!+\!1)$};

  % 水平键
  \draw[hbond] (TL) -- (TR);
  \draw[hbond] (BL) -- (BR);

  % 垂直键
  \draw[vbond] (TL) -- (BL);
  \draw[vbond] (TR) -- (BR);

  % "\"对角线键（红色）—— 从(i,j)到(i+1,j+1)
  \draw[dbond, very thick] (TL) -- (BR) node[pos=0.3, above right, red!70!black] {\small $\delta$};

  % "/"对角线键（橙色）—— 从(i,j+1)到(i+1,j)
  \draw[abond, very thick] (TR) -- (BL) node[pos=0.3, above left, tensororange!70!black] {\small $\gamma$};

  % 交叉点标记
  \fill[red!60] (0.5, 0.5) circle (0.08);
  \node[right, red!70!black] at (0.6, 0.5) {\textbf{交叉点}};
\end{tikzpicture}
\end{center}

$N \times N$ 棋盘有 $(N-1) \times (N-1)$ 个基本方块，因此有 $(N-1)^2$ 个键交叉。

\begin{center}
\begin{tabular}{c|ccccccc}
\toprule
$N$ & 3 & 4 & 5 & 8 & 10 & 16 & 20 \\
\midrule
交叉数 $(N-1)^2$ & 4 & 9 & 16 & 49 & 81 & 225 & 361 \\
\bottomrule
\end{tabular}
\end{center}

交叉数随 $N^2$ 增长——这是一个\textbf{体积量级}的非平面性，不是通过局部调整就能消除的。

\subsection{图论视角：Kuratowski 定理}

从图论的角度，$N$ 皇后张量网络的底层图是 $N \times N$ 方格子图与两组对角线边的并集。这个图是所谓的\textbf{国王图（king graph）}——每个顶点连接到它的 8 个邻居（上下左右 + 四个对角方向）。

由 Kuratowski 定理，一个图是平面的当且仅当它不包含 $K_5$ 或 $K_{3,3}$ 的细分作为子图。对于 $N \geq 3$ 的国王图，很容易找到 $K_{3,3}$ 的细分——例如取两行三列中的 6 个顶点及其对角线连接。因此 $N$ 皇后网络对所有 $N \geq 3$ 都是非平面的。

更进一步，这个图的\textbf{亏格（genus）}——即嵌入它所需的最小曲面的把手数——随 $(N-1)^2$ 增长（每个交叉至少需要一个把手来消除），是 $\Theta(N^2)$ 量级。


%=============================================================
\section{平面化策略：SWAP 张量}
%=============================================================

\subsection{核心思想}

既然非平面性来自对角线键的交叉，一个自然的解决方案是：在每个交叉点引入一个显式的\textbf{SWAP 张量}，使得两条键``立体交叉''变成``平面绕行''。

类比：在电路设计中，当两条导线在 PCB 上交叉时，可以用一个\textbf{过孔（via）}让一条导线穿过另一层。SWAP 张量在张量网络中扮演的角色与过孔相同。

\subsection{SWAP 张量的定义}

SWAP 张量 $S$ 是一个 4-指标张量，定义为：
\begin{equation}
\boxed{S_{\gamma \delta \gamma' \delta'} = \delta_{\gamma \gamma'} \cdot \delta_{\delta \delta'}}
\end{equation}
其中 $\gamma, \gamma' \in \{0,1\}$ 是 ``/'' 对角线的键指标，$\delta, \delta' \in \{0,1\}$ 是 ``$\backslash$'' 对角线的键指标。

\textbf{含义}：SWAP 张量不做任何计算——它只是让两条键各自独立地通过交叉点。$\gamma$ 和 $\gamma'$ 是同一条 ``/'' 键的输入和输出（必须相等），$\delta$ 和 $\delta'$ 是同一条 ``$\backslash$'' 键的输入和输出（也必须相等）。

SWAP 张量的大小为 $2 \times 2 \times 2 \times 2 = 16$ 个元素，其中只有 4 个非零：
\begin{equation}
S_{0000} = S_{0101} = S_{1010} = S_{1111} = 1, \quad \text{其余} = 0
\end{equation}

\begin{center}
\begin{tikzpicture}[scale=1.5]
  % 交叉点（SWAP张量）
  \node[swaptensor, minimum size=0.8cm] (S) at (0,0) {$S$};

  % "/" 对角线（橙色）穿过
  \draw[abond, ->] (-1,-1) -- (S.south west) node[near start, left] {$\gamma$};
  \draw[abond, ->] (S.north east) -- (1,1) node[near end, right] {$\gamma'$};

  % "\" 对角线（红色）穿过
  \draw[dbond, ->] (-1,1) -- (S.north west) node[near start, left] {$\delta$};
  \draw[dbond, ->] (S.south east) -- (1,-1) node[near end, right] {$\delta'$};

  % 等号
  \node at (2.5, 0) {\Large $=$};

  % 等价表示
  \draw[abond, ->] (3.5,-1) -- (4.5,0) -- (5.5,1);
  \draw[dbond, ->] (3.5,1) -- (4.5,0) -- (5.5,-1);
  \node[above] at (4.5, 0.3) {\small 各自独立穿过};
\end{tikzpicture}
\end{center}

\subsection{SWAP 张量的插入}

对于 $N \times N$ 棋盘，在每个基本方块的中心插入一个 SWAP 张量。具体操作：

\begin{enumerate}
\item 原来从 $(i,j)$ 直接连到 $(i+1,j+1)$ 的 ``$\backslash$'' 键，改为：$(i,j) \to S_{(i,j)} \to (i+1,j+1)$
\item 原来从 $(i,j+1)$ 直接连到 $(i+1,j)$ 的 ``/'' 键，改为：$(i,j+1) \to S_{(i,j)} \to (i+1,j)$
\item 其中 $S_{(i,j)}$ 是位于方块 $(i,j),(i,j+1),(i+1,j),(i+1,j+1)$ 中心的 SWAP 张量
\end{enumerate}

插入后，所有键都可以在平面内无交叉地画出——因为对角线键不再直接穿越，而是通过 SWAP 张量中转。

\begin{tcolorbox}[title=平面化的精确性, colback=green!5, colframe=darkgreen]
\textbf{关键性质}：SWAP 张量的插入不改变网络的收缩值。这是因为 $S_{\gamma \delta \gamma' \delta'} = \delta_{\gamma \gamma'} \delta_{\delta \delta'}$ 就是一个恒等操作——它在每条键上施加的约束（输入=输出）与不插入时完全一致。

数学上：对于任何张量网络，在一条键 $e$ 上插入恒等矩阵 $\delta_{ab}$ 不改变收缩值。SWAP 张量是两个独立恒等矩阵的张量积，因此同样不改变收缩值。
\end{tcolorbox}

\subsection{平面化网络的结构}

平面化后的网络具有\textbf{装饰方格子（decorated square lattice）}结构：

\begin{center}
\begin{tikzpicture}[scale=2.0]
  % 3x3格点张量
  \foreach \i in {0,1,2} {
    \foreach \j in {0,1,2} {
      \node[sitetensor, minimum size=0.5cm] (T\i\j) at (\j*2, -\i*2) {};
    }
  }

  % 2x2 SWAP张量（在方块中心）
  \foreach \i in {0,1} {
    \foreach \j in {0,1} {
      \node[swaptensor, minimum size=0.4cm] (S\i\j) at (\j*2+1, -\i*2-1) {};
    }
  }

  % 水平键
  \foreach \i in {0,1,2} {
    \foreach \j in {0,1} {
      \pgfmathtruncatemacro{\jj}{\j+1}
      \draw[hbond] (T\i\j) -- (T\i\jj);
    }
  }

  % 垂直键
  \foreach \j in {0,1,2} {
    \foreach \i in {0,1} {
      \pgfmathtruncatemacro{\ii}{\i+1}
      \draw[vbond] (T\i\j) -- (T\ii\j);
    }
  }

  % 对角线键（通过SWAP张量）
  \foreach \i in {0,1} {
    \foreach \j in {0,1} {
      \pgfmathtruncatemacro{\ii}{\i+1}
      \pgfmathtruncatemacro{\jj}{\j+1}
      % "\"方向
      \draw[dbond] (T\i\j) -- (S\i\j);
      \draw[dbond] (S\i\j) -- (T\ii\jj);
      % "/"方向
      \draw[abond] (T\i\jj) -- (S\i\j);
      \draw[abond] (S\i\j) -- (T\ii\j);
    }
  }

  \node[below] at (2, -5) {\small 平面化后的 $3 \times 3$ 网络：蓝色方块 = 格点张量 $\mathcal{T}$，橙色菱形 = SWAP 张量 $S$};
\end{tikzpicture}
\end{center}

\subsection{将 SWAP 张量融合到格点张量中}

为了使用标准的二维收缩算法，我们希望网络只有一种张量。一种方案是将 SWAP 张量\textbf{融合（absorb）}到邻近的格点张量中。

具体做法：将每个 SWAP 张量与其\textbf{右下方}的格点张量合并。设 SWAP 张量 $S_{(i,j)}$ 位于格点 $(i,j)$、$(i,j+1)$、$(i+1,j)$、$(i+1,j+1)$ 围成的方块中心。将它合并到格点 $(i+1,j+1)$ 的张量 $\mathcal{T}$ 中。

合并后，格点 $(i+1,j+1)$ 的新张量 $\widetilde{\mathcal{T}}$ 在接收对角线键的方向上需要更多指标来容纳来自 SWAP 的额外信息。具体来说：
\begin{itemize}
\item 原来格点 $(i+1,j+1)$ 的 ``$\backslash$'' 输入 $\delta_{\text{in}}$ 来自 $(i,j)$ 经过 $S_{(i,j)}$
\item 原来格点 $(i+1,j+1)$ 的 ``/'' 输入... 来自另一个 SWAP
\end{itemize}

实际上，更简洁的融合方案是：定义一个\textbf{超格点（super-site）}，将 $2 \times 2$ 的格点块（包括 4 个格点张量和 1 个 SWAP 张量）合并为一个大张量。这样：
\begin{itemize}
\item 超格点的物理指标：$2 \times 2 = 4$ 个二值变量
\item 对外键：只有水平和垂直方向（所有对角线键都在超格点内部被消化）
\item 结果：一个纯方格子网络，可直接使用 TRG/CTMRG
\end{itemize}

\begin{tcolorbox}[title=超格点方案, colback=yellow!5, colframe=tensororange]
\textbf{超格点张量 $\mathcal{T}_{\text{super}}$}：将 $2 \times 2$ 格点块中的 4 个格点张量 + 1 个 SWAP 张量完全收缩掉内部键，得到一个 8-腿张量（上下左右各两条）。

外部键维度：
\begin{itemize}
\item 水平方向：每侧有 2 条行约束键（来自两个格点），键维度 $2^2 = 4$
\item 垂直方向：每侧有 2 条列约束键 + 对角线信息，键维度 $2^2 \times 2 = 8$（或更小，取决于具体收缩顺序）
\end{itemize}

代价：有效键维度从 2 增大到 4--8，但网络变成了纯平面方格子。
\end{tcolorbox}


%=============================================================
\section{二维收缩方法}
%=============================================================

平面化之后，我们可以使用成熟的二维张量网络收缩方法。本节详细介绍三种主要方法。

\subsection{TRG（张量重正化群）}

TRG 由 Levin 和 Nave 在 2007 年提出，是最简单的二维张量网络粗粒化方法。

\textbf{基本思想}：在一个方格子张量网络中，将两个相邻张量通过 SVD 合并为一个更大的张量，然后截断到有限的键维度 $\chi$。重复此过程，网络的大小每步减半，直到缩小为一个可以直接计算的小网络。

\textbf{具体步骤}：

\textbf{第 1 步——SVD 分解}：对每个 4-腿张量 $T_{ijkl}$，沿一个方向做 SVD：
\begin{equation}
T_{ijkl} = \sum_{\alpha=1}^{\chi} U_{ij\alpha} \cdot V_{\alpha kl}
\end{equation}
这里 $U$ 和 $V$ 分别是``上半''和``下半''张量，$\alpha$ 是新的键指标。

\textbf{第 2 步——重组}：将相邻格点的``下半''和``上半''沿共享的键收缩：
\begin{equation}
T'_{i'j'k'l'} = \sum_{a,b} U_{i'a\alpha}(x,y) \cdot V_{\alpha aj'}(x,y+1) \cdot U_{k'b\beta}(x+1,y) \cdot V_{\beta bl'}(x+1,y+1)
\end{equation}
（这里省略了一些指标细节，只展示结构。）

\textbf{第 3 步——截断}：新张量 $T'$ 的键维度可能增大，通过 SVD 截断到 $\chi$。

\textbf{第 4 步——迭代}：重复步骤 1--3，每次迭代将网络的线性尺寸减半。对于 $N \times N$ 网络，需要 $\log_2 N$ 步。

\begin{tcolorbox}[title=TRG 的计算复杂度, colback=blue!5, colframe=tensorblue]
\begin{itemize}
\item 每步粗粒化的主要操作：张量收缩 + SVD 截断
\item 单步计算量：$O(\chi^6)$（主导项来自 SVD）
\item 总步数：$O(\log N)$
\item 总计算量：$O(\chi^6 \cdot \log N)$
\item 存储量：$O(\chi^4)$（一个 4-腿张量的大小）
\end{itemize}

对于平面化后的 $N$ 皇后网络，有效键维度 $D_{\text{eff}} \approx 4$--$8$。如果截断到 $\chi = 64$，单步计算量约 $64^6 \approx 7 \times 10^{10}$，在现代计算机上可以在秒级完成。
\end{tcolorbox}

\textbf{TRG 的局限性}：
\begin{enumerate}
\item \textbf{精度有限}：TRG 在 SVD 截断时没有考虑全局环境，只做了局部最优截断。对于关联较强的系统，误差累积较快。
\item \textbf{CDL 张量问题}：TRG 会产生所谓的 CDL（corner double line）结构，导致粗粒化后的张量中出现冗余自由度，降低精度。
\item \textbf{非均匀网络}：$N$ 皇后问题（有限 $N$）的网络是非均匀的（边界张量不同），TRG 需要跟踪每个张量的个性。
\end{enumerate}

\subsection{HOTRG（高阶张量重正化群）}

HOTRG 由 Xie 等人在 2012 年提出，是 TRG 的重要改进。

\textbf{关键改进}：HOTRG 不是先做 SVD 再重组，而是先将两个相邻张量完全收缩，然后用\textbf{高阶 SVD（HOSVD）}来截断。

\textbf{步骤}：
\begin{enumerate}
\item 选择一个方向（比如垂直方向），将上下相邻的两个张量收缩为一个 6-腿张量。
\item 对这个 6-腿张量做截断：通过计算其左右方向的约化密度矩阵，得到最优的截断投影算子 $P$。
\item 用 $P$ 将 6-腿张量压缩回 4-腿张量。
\item 在另一个方向重复。
\end{enumerate}

HOTRG 的截断考虑了张量的全部结构（通过约化密度矩阵），因此比 TRG 的纯 SVD 截断更优。

\textbf{计算复杂度}：$O(\chi^7)$ 每步，总共 $O(\chi^7 \log N)$。虽然比 TRG 多一个 $\chi$ 的因子，但精度的提升通常远超这一额外代价。

\subsection{CTMRG（角转移矩阵重正化群）}

CTMRG 最早由 Nishino 和 Okunishi 在 1996 年提出（基于 Baxter 1968 年的角转移矩阵思想），是目前最广泛使用的二维张量网络算法之一。它特别适合\textbf{热力学极限}下的计算，是研究皇后格点气体相变的理想工具。

\subsubsection{环境分解}

CTMRG 的核心思想是用有限数量的\textbf{环境张量}来近似无限大网络对中心区域的影响。这些环境张量包括：

\begin{itemize}
\item \textbf{四个角矩阵} $C_1, C_2, C_3, C_4$：二维张量（矩阵），大小为 $\chi \times \chi$
\item \textbf{四个边张量} $E_1, E_2, E_3, E_4$：三维张量，大小为 $\chi \times D \times \chi$
\end{itemize}

其中 $\chi$ 是环境键维度（可控参数），$D$ 是格点张量的物理键维度。

\begin{center}
\begin{tikzpicture}[scale=1.4]
  % 角矩阵
  \node[cornertensor] (C1) at (0,3) {$C_1$};
  \node[cornertensor] (C2) at (4,3) {$C_2$};
  \node[cornertensor] (C3) at (0,0) {$C_3$};
  \node[cornertensor] (C4) at (4,0) {$C_4$};

  % 边张量
  \node[edgetensor] (E1) at (2,3) {$E_1$};
  \node[edgetensor] (E2) at (4,1.5) {$E_2$};
  \node[edgetensor] (E3) at (2,0) {$E_3$};
  \node[edgetensor] (E4) at (0,1.5) {$E_4$};

  % 中心张量
  \node[sitetensor, minimum size=0.8cm] (T) at (2,1.5) {$\mathcal{T}$};

  % 连接
  \draw[thick] (C1) -- (E1) -- (C2);
  \draw[thick] (C2) -- (E2) -- (C4);
  \draw[thick] (C4) -- (E3) -- (C3);
  \draw[thick] (C3) -- (E4) -- (C1);
  \draw[thick] (E1) -- (T) -- (E3);
  \draw[thick] (E4) -- (T) -- (E2);

  \node[below] at (2, -0.7) {CTMRG 环境结构：角矩阵 $C_k$ + 边张量 $E_k$ 包围中心格点 $\mathcal{T}$};
\end{tikzpicture}
\end{center}

\subsubsection{迭代算法}

CTMRG 通过\textbf{吸收（absorption）}步骤逐渐扩大被近似的区域。以向右扩展一列为例：

\textbf{第 1 步——吸收}：将新的一列格点张量 $\mathcal{T}$ 合并到环境中：
\begin{align}
\widetilde{C}_1 &= C_1 \cdot E_1 \cdot \mathcal{T} \quad (\text{角矩阵吸收边张量和格点张量}) \\
\widetilde{E}_4 &= E_4 \cdot \mathcal{T} \cdot E_2 \quad (\text{边张量吸收格点张量})
\end{align}
（这里用简化记号表示，实际涉及多个指标的收缩。）

\textbf{第 2 步——截断}：吸收后，环境键维度从 $\chi$ 增大到 $\chi \cdot D$。通过 SVD 截断将其压缩回 $\chi$：
\begin{enumerate}
\item 构造密度矩阵 $\rho = \widetilde{C}_1 \widetilde{C}_1^\dagger + \widetilde{C}_4 \widetilde{C}_4^\dagger$
\item 对 $\rho$ 做特征值分解，保留最大的 $\chi$ 个特征向量作为投影算子 $P$
\item $C_1^{\text{new}} = P^\dagger \widetilde{C}_1 P$，$E_4^{\text{new}} = P^\dagger \widetilde{E}_4 P$
\end{enumerate}

\textbf{第 3 步——归一化}：将更新后的角矩阵和边张量归一化，防止数值溢出。

\textbf{第 4 步——重复}：对四个方向轮流执行吸收-截断步骤，直到收敛。

\begin{tcolorbox}[title=CTMRG 收敛判据, colback=green!5, colframe=darkgreen]
常用的收敛判据：
\begin{itemize}
\item 角矩阵的奇异值谱不再变化：$||\sigma^{(n)} - \sigma^{(n-1)}|| < \epsilon$
\item 物理观测量（如密度 $\rho$）的变化小于阈值
\item 典型阈值：$\epsilon \sim 10^{-8}$--$10^{-12}$
\end{itemize}
\end{tcolorbox}

\subsubsection{计算复杂度}

\begin{center}
\begin{tabular}{ll}
\toprule
操作 & 复杂度 \\
\midrule
吸收步骤 & $O(\chi^2 D^2 + \chi^2 D^3)$ \\
密度矩阵构造 & $O(\chi^2 D)$ \\
SVD 截断 & $O(\chi^3 D^3)$ \\
每轮迭代（4 方向） & $O(\chi^3 D^3)$ \\
\bottomrule
\end{tabular}
\end{center}

对于平面化后的皇后格点张量，$D \approx 4$--$8$。如果取 $\chi = 100$，每轮迭代的计算量约 $10^6 \times 8^3 \approx 5 \times 10^8$，非常快。

\subsubsection{适配 N 皇后问题的注意事项}

\textbf{1. 非均匀性处理}：
对于有限 $N$ 的精确计数，每个格点张量、每条边界向量都不同（取决于位置）。这时需要使用\textbf{非均匀 CTMRG}：为每个格点维护独立的环境张量。这将存储量从 $O(\chi^2)$ 增加到 $O(N^2 \chi^2)$。

对于皇后格点气体（$N \to \infty$，所有约束放宽为``至多一个''），网络是\textbf{均匀的}（平移不变），只需维护一套环境张量即可。

\textbf{2. 超格点的使用}：
如第 4 节讨论的，将 $2 \times 2$ 格点块合并为超格点后，网络变为纯方格子。CTMRG 可以直接在这个超格子上运行，无需任何修改。超格点的有效键维度决定了 $D$ 的大小。

\textbf{3. 对称性利用}（详见第 9 节）：
如果利用棋盘的 $D_4$ 对称性，可以令 $C_1 = C_2 = C_3 = C_4$ 和 $E_1 = E_2 = E_3 = E_4$，将计算量和存储量减少 4--8 倍。

\subsection{三种方法的对比}

\begin{center}
\begin{tabular}{lccc}
\toprule
 & TRG & HOTRG & CTMRG \\
\midrule
精度（同 $\chi$） & 较低 & 中等 & 较高 \\
单步复杂度 & $O(\chi^6)$ & $O(\chi^7)$ & $O(\chi^3 D^3)$ \\
适用场景 & 有限系统 & 有限系统 & 无限系统 / 热力学极限 \\
非均匀网络 & 自然支持 & 自然支持 & 需要修改 \\
实现难度 & 简单 & 中等 & 中等 \\
CDL 问题 & 有 & 有（较轻） & 无 \\
\bottomrule
\end{tabular}
\end{center}

\textbf{推荐}：
\begin{itemize}
\item 有限 $N$ 的近似计数：首选 HOTRG
\item 皇后格点气体的热力学性质：首选 CTMRG
\item 快速原型验证：首选 TRG（最简单）
\end{itemize}


%=============================================================
\section{皇后格点气体：热力学极限与相变}
%=============================================================

\subsection{模型定义回顾}

$N$ 皇后问题要求恰好放 $N$ 个皇后。皇后格点气体是其自然推广：在 $N \times N$ 棋盘上放\textbf{任意数量}的互不攻击皇后，每个皇后携带逸度 $z$。配分函数为：
\begin{equation}
\boxed{\mathcal{Z}(z) = \sum_{k=0}^{N} z^k \cdot Q(N, k)}
\end{equation}
其中 $Q(N, k)$ 是放 $k$ 个互不攻击皇后的方案数。特别地，$Q(N, N) = Z_N$（$N$ 皇后解的总数）。

使用格点变量，配分函数可以写为：
\begin{equation}
\mathcal{Z}(z) = \sum_{\{n_{ij}\}} z^{\sum_{ij} n_{ij}} \prod_{\text{约束}} \delta(\text{无冲突})
\end{equation}

\textbf{物理图像}：这是一个\textbf{硬核格点气体}——皇后之间的互不攻击约束等价于无穷大的排斥势（硬核势）。逸度 $z$ 控制系统的密度：$z \to 0$ 时棋盘几乎为空，$z \to \infty$ 时系统趋向最密堆积。

\subsection{张量网络表示}

格点气体的张量网络与 $N$ 皇后问题几乎完全相同，只需两处修改：

\textbf{1. 约束放宽}：行约束和列约束从``恰好一个''改为``至多一个''。这只需修改 MPS 的边界向量：
\begin{equation}
|R\rangle = \begin{pmatrix} 0 \\ 1 \end{pmatrix} \quad \longrightarrow \quad |R'\rangle = \begin{pmatrix} 1 \\ 1 \end{pmatrix}
\end{equation}
（接受最终状态为 0 或 1——即该行/列可以没有皇后。）

\textbf{2. 逸度权重}：每放一个皇后贡献因子 $z$。将 $z$ 均匀分配到四条约束链上，每条贡献 $z^{1/4}$：
\begin{equation}
A_z^{[0]} = \begin{pmatrix} 1 & 0 \\ 0 & 1 \end{pmatrix}, \quad
A_z^{[1]} = \begin{pmatrix} 0 & z^{1/4} \\ 0 & 0 \end{pmatrix}
\end{equation}
每个皇后的总权重：$(z^{1/4})^4 = z$。

格点气体的格点张量：
\begin{equation}
\mathcal{T}_z = \sum_{n=0,1} (A_z^{[n]})_{\text{row}} \cdot (A_z^{[n]})_{\text{col}} \cdot (B_z^{[n]})_{\text{diag}/} \cdot (B_z^{[n]})_{\text{diag}\backslash}
\end{equation}

\textbf{关键特征}：由于所有约束都是``至多一个''，且所有格点的张量完全相同，网络具有\textbf{完美的平移不变性}——这正是 CTMRG 大展身手的场景。

\subsection{相变行为}

皇后格点气体在热力学极限 $N \to \infty$ 下展现丰富的相变行为。

\subsubsection{低密度相（$z < z_c$）}

当 $z$ 较小时，皇后密度低，互不攻击的约束不太紧张。系统类似于普通的稀薄格点气体：
\begin{itemize}
\item 皇后在棋盘上近似均匀分布
\item 关联函数 $G(r) = \langle n_0 n_r \rangle - \langle n_0 \rangle \langle n_r \rangle$ 呈\textbf{指数衰减}：$G(r) \sim e^{-r/\xi}$
\item 关联长度 $\xi(z)$ 有限
\item 无长程序
\end{itemize}

\subsubsection{高密度相（$z > z_c$）}

当 $z$ 超过临界值 $z_c$ 时，皇后密度高到系统被迫形成有序结构：
\begin{itemize}
\item 皇后排列形成\textbf{有序模式}（类似于阶梯形排列）
\item 平移对称性被\textbf{自发破缺}
\item 关联函数呈长程序：$G(r) \to \text{const} > 0$（$r \to \infty$）
\item 存在非零\textbf{序参量}
\end{itemize}

\textbf{物理直觉}：皇后的攻击范围覆盖了整行、整列和整条对角线，这意味着每放一个皇后就``封锁''了很大一片区域。在高密度下，皇后必须以极其有序的方式排列才能互不攻击——这类似于\textbf{熵驱动的有序化}（entropy-driven ordering），在硬核粒子系统中普遍存在。

\subsubsection{临界点 $z_c$}

临界点 $z_c$ 处，系统发生连续相变：
\begin{itemize}
\item 关联长度 $\xi \to \infty$
\item 序参量从 0 连续增长
\item 比热出现峰值或发散
\item 系统处于\textbf{无标度}状态
\end{itemize}

$z_c$ 的精确值目前未知，这正是张量网络方法（CTMRG）可以提供新信息的地方。

\subsection{CTMRG 可以计算什么？}

使用 CTMRG 方法，我们可以直接计算以下热力学量：

\textbf{1. 每格点自由能}：
\begin{equation}
f(z) = -\lim_{N \to \infty} \frac{1}{N^2} \ln \mathcal{Z}(z)
\end{equation}
通过 CTMRG 环境张量的归一化因子直接获得。

\textbf{2. 皇后密度}：
\begin{equation}
\rho(z) = \frac{z}{\mathcal{Z}} \frac{\partial \mathcal{Z}}{\partial z} = -z \frac{\partial f}{\partial z}
\end{equation}
通过在网络中央插入一个``杂质张量''（$n_{ij}$ 的期望值）计算。

\textbf{3. 涨落（压缩率/比热的类比）}：
\begin{equation}
\kappa = \frac{\partial \rho}{\partial (\ln z)} = \langle k^2 \rangle - \langle k \rangle^2 \propto \sum_r G(r)
\end{equation}

\textbf{4. 关联长度}：
\begin{equation}
\xi = -\frac{1}{\ln(\lambda_2 / \lambda_1)}
\end{equation}
其中 $\lambda_1, \lambda_2$ 是转移矩阵的前两个特征值。CTMRG 的边张量 $E$ 的谱直接给出这个信息。

\textbf{5. 序参量}：定义子格磁化率或类似量来探测对称性破缺。

\subsection{与蒙特卡洛的互补性}

皇后格点气体在高密度区域的蒙特卡洛模拟面临严重的\textbf{遍历性问题}：合法构型之间被大量不合法构型隔开，单步翻转（添加/删除一个皇后）几乎无法接受。这导致：
\begin{itemize}
\item 自关联时间 $\tau$ 在 $z > z_c$ 时指数增长
\item 模拟结果严重受限于初始构型的选择
\item 即使使用高级算法（蠕虫算法、集团更新），在高密度区域仍然困难
\end{itemize}

张量网络方法\textbf{完全绕过了遍历性问题}，因为它直接计算配分函数，不需要在构型空间中采样。这使得张量网络成为研究高密度相的唯一可靠工具。


%=============================================================
\section{纠缠结构与计算复杂度}
%=============================================================

本节从信息论和计算复杂度的角度深入分析 $N$ 皇后问题。

\subsection{$N$ 皇后解空间的``纠缠熵''}

定义``$N$ 皇后态''为所有合法构型的等权叠加：
\begin{equation}
|\Psi_N\rangle = \frac{1}{\sqrt{Z_N}} \sum_{\text{合法构型 } \sigma} |\sigma\rangle
\end{equation}
其中 $|\sigma\rangle = |n_{11} n_{12} \cdots n_{NN}\rangle$ 是一个 $N^2$ 量子比特态。

考虑将棋盘沿垂直中线（列 $N/2$ 和 $N/2+1$ 之间）切开，得到左半 $L$ 和右半 $R$。``纠缠熵''为：
\begin{equation}
S(L:R) = -\sum_i p_i \ln p_i
\end{equation}
其中 $\{p_i\}$ 是约化密度矩阵 $\rho_L = \text{Tr}_R |\Psi_N\rangle\langle\Psi_N|$ 的特征值。

\textbf{Schmidt 秩的下界}：纠缠熵至少为 Schmidt 秩的对数。Schmidt 秩等于切割处需要跟踪的不同``边界模式''的数目。

在 $N$ 皇后问题中，切割处需要传递的信息至少包括：哪些列已经被左半部分的皇后占据？由于前 $N/2$ 列中恰好有某些行的皇后占据了这些列，不同的列占据模式数为 $\binom{N}{k}$（$k$ 依赖于具体的行分布，但量级为 $\binom{N}{N/2}$）。因此：
\begin{equation}
\boxed{S(L:R) \geq \ln \binom{N}{N/2} \sim N \ln 2 - \frac{1}{2} \ln(\pi N / 2)}
\end{equation}

这是\textbf{线性增长}的——与系统大小 $N$ 成正比。这意味着 $N$ 皇后态具有\textbf{体积律纠缠}（在一维切割的意义下），与物理基态的面积律纠缠形成鲜明对比。

\subsection{与 \#P-完全性的联系}

\#P 是一个计算复杂度类，包含所有``计算 NP 问题的解的个数''的问题。$\#N$-Queens（计算 $N$ 皇后解的个数）是 \#P-完全的。

\begin{tcolorbox}[title=\#P-完全性与张量网络, colback=red!5, colframe=red!60!black]
\textbf{定理}（非形式化）：如果存在一个张量网络方法，能在多项式时间内精确计算 $Z_N$（使用多项式大小的键维度 $\chi = \text{poly}(N)$），那么 \#P = FP（即所有 \#P 问题都可以在多项式时间内求解）。

\textbf{推论}：假设 \#P $\neq$ FP（这是一个被广泛相信但未证明的假设），则\textbf{精确}计算 $Z_N$ 的任何张量网络方法都必须使用指数大的键维度。

\textbf{重要注意}：这个障碍只限于\textbf{精确}计数。对于\textbf{近似}计算——特别是热力学量（自由能、密度等每格点量）——\textbf{不存在这样的复杂度障碍}。事实上，许多 \#P-硬的统计力学模型（如反铁磁 Potts 模型）的自由能可以用 $\chi = O(1)$ 的 CTMRG 很好地近似。
\end{tcolorbox}

\subsection{纠缠谱分析}

纠缠谱（entanglement spectrum）——即 Schmidt 值的分布——揭示了问题的内在结构。

对于物理系统（如 Ising 模型临界点），纠缠谱通常呈现：
\begin{itemize}
\item 少数几个大的 Schmidt 值（主导贡献）
\item 快速衰减的尾部
\item 明显的``纠缠隙''（entanglement gap）
\end{itemize}
这种结构使得 MPS/PEPS 截断非常有效——丢弃小 Schmidt 值引入的误差很小。

对于 $N$ 皇后问题，我们预期纠缠谱更加``平坦''：
\begin{itemize}
\item 许多 Schmidt 值具有可比的大小
\item 没有明显的纠缠隙
\item 截断任何 Schmidt 值都会引入显著误差
\end{itemize}

这反映了 $N$ 皇后问题的组合本质——每种列占据模式（每种``边界条件''）对应的解的数目不存在明显的层次结构。

\subsection{近似计算的可能性}

尽管精确计数需要指数资源，近似计算可能需要的资源要少得多。关键区分在于\textbf{我们要计算什么}：

\textbf{精确计数 $Z_N$}：需要 $\chi \sim 2^N$，不可避免。

\textbf{$\ln Z_N$ 的近似}：$\ln Z_N$ 是一个$O(N \ln N)$ 量级的数，即使相对误差为 $10^{-3}$，对应的 $Z_N$ 的相对误差也可能很大。但如果我们只关心 $\ln Z_N / N^2$（每格点自由能），所需的 $\chi$ 可能远小于 $2^N$。

\textbf{皇后密度 $\rho(z)$}：作为一个\textbf{局部}物理量的期望值，$\rho(z)$ 对远处的信息不敏感。CTMRG 用有限的 $\chi$ 就能给出良好的近似——这已经在无数个统计力学模型中得到验证。

\textbf{临界指数}：临界行为由长波长涨落主导，CTMRG 的有限 $\chi$ 引入了一个有效的红外截断。通过 $\chi$ 外推可以得到临界指数的估计。


%=============================================================
\section{与蒙特卡洛方法的系统对比}
%=============================================================

\subsection{蒙特卡洛方法的简要回顾}

蒙特卡洛（MC）方法通过随机抽样来估计统计力学量。对于皇后格点气体，最基本的 Metropolis 算法如下：

\begin{enumerate}
\item 从一个合法构型出发
\item 随机选择一个格点 $(i,j)$
\item 如果 $(i,j)$ 上有皇后，提议移除：以概率 $\min(1, 1/z)$ 接受
\item 如果 $(i,j)$ 上无皇后且不与任何现有皇后冲突，提议放置：以概率 $\min(1, z)$ 接受
\item 重复步骤 2--4，收集统计量
\end{enumerate}

\subsection{遍历性问题的深入分析}

\textbf{低密度}（$z < z_c$）：算法运行良好。皇后稀疏，每次添加/删除皇后的接受率较高。

\textbf{高密度}（$z > z_c$）：算法陷入困境。原因如下：

\begin{enumerate}
\item \textbf{几何受挫}：每个皇后封锁了整行、整列和两条对角线。在高密度下，几乎所有空位都被至少一个皇后攻击，无法添加新皇后。
\item \textbf{删除代价高}：删除一个皇后的 Boltzmann 权重为 $1/z$，当 $z \gg 1$ 时，删除几乎不被接受。
\item \textbf{结构刚性}：合法的高密度构型形成了高度有序的结构（类似于水晶），从一种有序结构跃迁到另一种需要大规模的协同重排——这在单步翻转的 MC 中几乎不可能。
\end{enumerate}

量化地说，自关联时间 $\tau(z)$ 在 $z > z_c$ 时预期呈\textbf{指数增长}：
\begin{equation}
\tau(z) \sim e^{c N} \quad (z > z_c)
\end{equation}
这使得 MC 模拟在热力学极限下对高密度相完全失效。

\subsection{高级蒙特卡洛方法}

\textbf{蠕虫算法（Worm Algorithm）}：引入一对``缺陷''（一个多余皇后和一个空穴），允许临时的约束违反。缺陷可以在棋盘上``蠕动''，连接不同的合法构型。这可以部分缓解遍历性问题，但对于皇后问题的长程约束效果有限。

\textbf{集团更新}：类似于 Ising 模型的 Wolff 算法，一次翻转一大群皇后。但皇后问题的约束结构使得定义合适的集团非常困难。

\textbf{并行回火（Parallel Tempering）}：同时模拟多个不同 $z$ 值的系统，允许构型在系统间交换。这可以帮助高 $z$ 的系统借用低 $z$ 系统的遍历性。

\textbf{Wang-Landau 采样}：直接估计态密度 $Q(N,k)$，通过平坦直方图判据调节采样权重。这可以避免高密度区域的采样困难，但收敛速度随系统大小变慢。

\subsection{系统对比}

\begin{center}
\begin{tabular}{p{3.5cm}p{5cm}p{5cm}}
\toprule
 & \textbf{蒙特卡洛} & \textbf{张量网络} \\
\midrule
基本原理 & 随机采样 & 确定性收缩 \\
误差类型 & 统计误差 $\sim 1/\sqrt{N_{\text{samples}}}$ & 系统误差 $\sim f(\chi)$ \\
精确计数 & 不可能 & 可能（小 $N$） \\
热力学极限 & 间接（有限尺寸标度） & 直接（CTMRG） \\
高密度相 & 遍历性障碍 & 无此问题 \\
临界区域 & 临界慢化 & $\chi$ 外推 \\
编程难度 & 简单 & 中等 \\
计算资源 & CPU 时间 & CPU 时间 + 内存 \\
可推广性 & 高（任何统计模型） & 中（需要张量网络构造） \\
\bottomrule
\end{tabular}
\end{center}

\textbf{结论}：对于皇后格点气体的研究，张量网络方法和蒙特卡洛方法是\textbf{互补的}。在低密度和临界区域，两者都可以工作；在高密度区域，张量网络是唯一可靠的方法。最佳策略是：在两种方法都适用的区域进行交叉验证，然后用张量网络深入探索 MC 无法到达的区域。


%=============================================================
\section{对称性的利用——$D_4$ 群}
%=============================================================

\subsection{棋盘的对称群}

$N \times N$ 正方棋盘具有二面体群 $D_4$ 对称性，包含 8 个元素：

\begin{center}
\begin{tabular}{cl}
\toprule
元素 & 操作 \\
\midrule
$e$ & 恒等变换 \\
$r$ & 逆时针旋转 $90°$ \\
$r^2$ & 旋转 $180°$ \\
$r^3$ & 逆时针旋转 $270°$（= 顺时针 $90°$） \\
$s$ & 水平轴反射 \\
$sr$ & 右对角线反射 \\
$sr^2$ & 垂直轴反射 \\
$sr^3$ & 左对角线反射 \\
\bottomrule
\end{tabular}
\end{center}

$N$ 皇后问题的解空间在 $D_4$ 作用下封闭：如果 $\sigma$ 是一个合法摆法，那么 $g\sigma$（对 $\sigma$ 施加任何 $D_4$ 变换 $g$）也是合法摆法。

\subsection{对称性在张量网络中的体现}

格点张量 $\mathcal{T}$ 在 $D_4$ 变换下具有确定的变换规则。以 $90°$ 旋转 $r$ 为例：
\begin{equation}
r: \quad (\alpha_L, \alpha_R, \beta_U, \beta_D, \gamma_{\text{in}}, \gamma_{\text{out}}, \delta_{\text{in}}, \delta_{\text{out}}) \mapsto (\beta_U, \beta_D, \alpha_R, \alpha_L, \delta_{\text{in}}, \delta_{\text{out}}, \gamma_{\text{out}}, \gamma_{\text{in}})
\end{equation}
（水平键变垂直键，``/'' 对角线变 ``$\backslash$'' 对角线，等等。注意方向的变化。）

由于格点张量的对称性（$A$ 和 $B$ 矩阵形式相同），$\mathcal{T}$ 在 $D_4$ 下是不变的（在适当重新标记指标后）。

\subsection{对称性在 CTMRG 中的利用}

对于具有 $D_4$ 对称性的均匀网络（皇后格点气体），CTMRG 的环境张量满足：
\begin{equation}
C_1 = C_2 = C_3 = C_4 \equiv C, \quad E_1 = E_2 = E_3 = E_4 \equiv E
\end{equation}

这意味着：
\begin{itemize}
\item 只需存储和更新\textbf{一个}角矩阵 $C$ 和\textbf{一个}边张量 $E$
\item 每轮迭代只需执行一个方向的吸收-截断
\item 计算量和存储量减少约 4 倍
\end{itemize}

更进一步，反射对称性要求 $C$ 是对称矩阵：$C = C^T$。这可以在 SVD 截断时强制施加，提高数值稳定性。

\subsection{$\mathbb{Z}_2$ 对称性与块对角化}

格点张量还具有一个内部 $\mathbb{Z}_2$ 对称性：将所有键指标同时翻转（$0 \leftrightarrow 1$）等价于交换``有皇后''和``无皇后''的角色。更准确地说，由于每个约束链跟踪``是否已经见过皇后''，状态 0 和 1 之间存在一种单调性结构。

虽然 $\mathbb{Z}_2$ 对称性不能直接将张量块对角化（因为 $n = 0$ 和 $n = 1$ 的贡献耦合在一起），但约束的单调性（状态只能从 0 跳到 1，不能反向）意味着转移矩阵是\textbf{上三角}的在某种意义上——这可以用来优化计算。

\subsection{Burnside 引理：不等价解的计数}

利用 $D_4$ 对称性，可以计算本质不同的 $N$ 皇后解的数目（不考虑旋转和反射的等价性）：
\begin{equation}
Z_N^{\text{distinct}} = \frac{1}{8}\left(\sum_{g \in D_4} |\text{Fix}(g)|\right)
\end{equation}
其中 $|\text{Fix}(g)|$ 是在变换 $g$ 下不动的解的数目。

对于大多数 $N$，$|\text{Fix}(g)| = 0$（当 $g \neq e$），此时 $Z_N^{\text{distinct}} = Z_N / 8$。但某些 $N$ 可能存在具有额外对称性的特殊解。

张量网络方法可以通过在收缩过程中\textbf{投影到 $D_4$ 不可约表示}来分别计算每种对称性扇区的贡献。


%=============================================================
\section{四种具体计算方案}
%=============================================================

\subsection{方案一：精确计数——MPS-MPO 方法}

这是最直接的方法，适用于小到中等的 $N$（$N \leq 20$ 左右）。

\begin{algorithm}[H]
\caption{精确 MPS-MPO 计算 $Z_N$}
\begin{algorithmic}[1]
\REQUIRE 棋盘大小 $N$
\ENSURE $N$ 皇后解的总数 $Z_N$
\STATE 构造 MPO 局部张量 $W^{[j]}$（$j = 1, \ldots, N$），键维度 $D_{\text{MPO}} = 4$
\STATE 初始化 MPS：$|\psi_0\rangle = |\mathbf{0}\rangle$（所有列空闲，无对角线威胁），$\chi_0 = 1$
\FOR{$m = 1$ to $N$}
    \STATE $|\widetilde{\psi}_m\rangle = T |\psi_{m-1}\rangle$ \COMMENT{MPS-MPO 乘法}
    \STATE 新键维度 $\chi_m = \chi_{m-1} \times D_{\text{MPO}} \times d$（$d = 8$ 为局部希尔伯特空间维度）
    \STATE 对 $|\widetilde{\psi}_m\rangle$ 做 SVD 正则化（精确，不截断）
    \STATE $|\psi_m\rangle = |\widetilde{\psi}_m\rangle$
\ENDFOR
\STATE $Z_N = \langle \text{final} | \psi_N \rangle$
\RETURN $Z_N$
\end{algorithmic}
\end{algorithm}

\textbf{关键实现细节}：
\begin{itemize}
\item MPS-MPO 乘法的核心操作：对每个位点 $j$，将 MPS 张量 $M^{[j]}$（维度 $\chi \times d \times \chi$）与 MPO 张量 $W^{[j]}$（维度 $D \times d \times d \times D$）收缩，得到新的 MPS 张量（维度 $\chi D \times d \times \chi D$）。
\item 稀疏操作：格点张量只有 17 个非零元素，可以利用稀疏张量运算大幅减少计算量。
\item 对称性：利用列占据向量的排列对称性可以减少独立计算量。
\end{itemize}

\textbf{预期性能}：
\begin{center}
\begin{tabular}{c|cccccc}
\toprule
$N$ & 8 & 10 & 12 & 14 & 16 & 18 \\
\midrule
$\chi_{\max}$ & 568 & $\sim 10^3$ & $\sim 10^4$ & $\sim 10^5$ & $\sim 10^6$ & $\sim 10^7$ \\
预估时间 & $<$ 1秒 & 秒 & 分钟 & 小时 & 天 & 周 \\
\bottomrule
\end{tabular}
\end{center}

\subsection{方案二：近似计数——截断 MPS-MPO}

在方案一的基础上增加 SVD 截断：

\begin{algorithm}[H]
\caption{近似 MPS-MPO 计算 $Z_N$}
\begin{algorithmic}[1]
\REQUIRE 棋盘大小 $N$，最大键维度 $\chi_{\max}$
\ENSURE $Z_N$ 的近似值
\STATE 同方案一初始化
\FOR{$m = 1$ to $N$}
    \STATE $|\widetilde{\psi}_m\rangle = T |\psi_{m-1}\rangle$ \COMMENT{MPS-MPO 乘法}
    \IF{$\chi > \chi_{\max}$}
        \STATE 对 $|\widetilde{\psi}_m\rangle$ 做 SVD，保留前 $\chi_{\max}$ 个奇异值
    \ENDIF
    \STATE $|\psi_m\rangle = |\widetilde{\psi}_m\rangle$
\ENDFOR
\STATE $Z_N \approx \langle \text{final} | \psi_N \rangle$
\RETURN $Z_N$
\end{algorithmic}
\end{algorithm}

\textbf{误差控制}：
\begin{itemize}
\item 运行 $\chi_{\max} = \chi$ 和 $\chi_{\max} = 2\chi$ 两次，比较结果
\item 如果相对差 $|Z(2\chi) - Z(\chi)| / |Z(2\chi)| < \epsilon$，认为结果可靠
\item 典型经验：$\chi_{\max} = 256$ 对 $N \leq 20$ 可以给出 $\sim 1\%$ 的相对误差
\end{itemize}

\subsection{方案三：皇后格点气体——CTMRG}

\begin{algorithm}[H]
\caption{CTMRG 计算皇后格点气体}
\begin{algorithmic}[1]
\REQUIRE 逸度 $z$，环境键维度 $\chi$，收敛阈值 $\epsilon$
\ENSURE 皇后密度 $\rho(z)$
\STATE 构造格点张量 $\mathcal{T}_z$（含逸度权重 $z^{1/4}$）
\STATE 如果使用超格点方案，合并 $2 \times 2$ 块消除对角线
\STATE 初始化：$C = \mathcal{T}_z$ 的适当收缩，$E = \mathcal{T}_z$ 的适当收缩
\REPEAT
    \FOR{方向 $d$ = 上, 右, 下, 左}
        \STATE 吸收：将新的 $\mathcal{T}_z$ 合并到 $C$ 和 $E$ 中
        \STATE 构造密度矩阵 $\rho_d$
        \STATE SVD 截断到 $\chi$：得到投影算子 $P_d$
        \STATE 更新：$C \leftarrow P_d^\dagger \widetilde{C} P_d$，$E \leftarrow P_d^\dagger \widetilde{E} P_d$
        \STATE 归一化 $C$ 和 $E$
    \ENDFOR
    \STATE 检查收敛：$C$ 的奇异值谱变化 $< \epsilon$？
\UNTIL{收敛}
\STATE 计算 $\rho = \text{Tr}(\mathcal{T}_z^{\text{imp}} \cdot \text{环境}) / \text{Tr}(\mathcal{T}_z \cdot \text{环境})$
\RETURN $\rho(z)$
\end{algorithmic}
\end{algorithm}

其中 $\mathcal{T}_z^{\text{imp}}$ 是``杂质张量''，与 $\mathcal{T}_z$ 相同但额外乘以格点上的观测量（如 $n_{ij}$）。

\textbf{关键参数选择}：
\begin{center}
\begin{tabular}{lll}
\toprule
参数 & 典型值 & 说明 \\
\midrule
$\chi$ & 32--256 & 越大越精确，计算量 $\sim \chi^3$ \\
$\epsilon$ & $10^{-8}$--$10^{-12}$ & 收敛阈值 \\
$z$ 范围 & $0.01$--$100$ & 覆盖低密度到高密度 \\
$z$ 间距 & 0.01（临界附近更密） & 分辨相变 \\
\bottomrule
\end{tabular}
\end{center}

\subsection{方案四：平面化网络 + TRG}

\begin{algorithm}[H]
\caption{平面化 + TRG 计算 $Z_N$}
\begin{algorithmic}[1]
\REQUIRE 棋盘大小 $N$，截断维度 $\chi$
\ENSURE $Z_N$ 的近似值
\STATE 构造 $N^2$ 个格点张量 $\mathcal{T}_{(i,j)}$ + $(N-1)^2$ 个 SWAP 张量 $S$
\STATE 将 SWAP 张量融合到格点张量中（或采用超格点方案）
\STATE 添加边界向量
\WHILE{网络中张量数 $> 1$}
    \STATE 选择收缩方向（交替水平/垂直）
    \STATE 对每对相邻张量：SVD 分解 $\to$ 截断到 $\chi$ $\to$ 重组
    \STATE 网络大小减半
\ENDWHILE
\RETURN 最终标量值 $\approx Z_N$
\end{algorithmic}
\end{algorithm}

\textbf{注意}：对于有限 $N$ 的非均匀网络，TRG 的每个张量都不同，需要分别跟踪。这增加了实现复杂度和存储量，但算法框架不变。

\subsection{四种方案的总结对比}

\begin{center}
\begin{tabular}{p{2.5cm}cccc}
\toprule
 & 方案一 & 方案二 & 方案三 & 方案四 \\
 & MPS-MPO精确 & MPS-MPO截断 & CTMRG & TRG \\
\midrule
目标 & 精确 $Z_N$ & 近似 $Z_N$ & $\rho(z)$, $f(z)$ & 近似 $Z_N$ \\
适用范围 & $N \leq 18$ & $N \leq 30+$ & $N \to \infty$ & $N \leq 30+$ \\
维度 & 1D & 1D & 2D & 2D \\
关键参数 & 无 & $\chi_{\max}$ & $\chi$ & $\chi$ \\
时间复杂度 & $O(N \cdot 2^{3N})$ & $O(N^2 \chi^3)$ & $O(\chi^3 D^3)$ & $O(\chi^6 \log N)$ \\
空间复杂度 & $O(N \cdot 2^N)$ & $O(N \chi)$ & $O(\chi^2)$ & $O(\chi^4)$ \\
实现难度 & 简单 & 简单 & 中等 & 中等 \\
\bottomrule
\end{tabular}
\end{center}


%=============================================================
\section{讨论与展望}
%=============================================================

\subsection{关键洞察总结}

本文的核心洞察可以概括为以下几点：

\begin{enumerate}
\item \textbf{非平面性是根本障碍}：$N$ 皇后张量网络因对角线交叉而是非平面的，$(N-1)^2$ 个交叉点使得标准二维方法不能直接使用。这一点在文献中很少被讨论。

\item \textbf{平面化是可行的}：通过引入 SWAP 张量，可以将非平面网络精确地转化为平面网络。代价是有效键维度的增大（从 2 增大到 4--8），但不改变收缩值。

\item \textbf{二维方法有独特优势}：特别是对皇后格点气体的热力学极限，CTMRG 可以直接计算每格点量，绕过了蒙特卡洛的遍历性问题。

\item \textbf{精确计数的指数困难是本质的}：\#P-完全性保证了任何方法都需要指数资源。但近似计算——特别是热力学量——可能只需要多项式资源。

\item \textbf{纠缠结构揭示问题本质}：$N$ 皇后问题的体积律纠缠（$S \sim N$）反映了其组合复杂性，与物理系统的面积律形成对比。
\end{enumerate}

\subsection{开放问题}

\textbf{1. 皇后格点气体的普适性类（universality class）}：
临界点 $z_c$ 附近的临界指数是什么？由于皇后的长程攻击范围，这个系统可能不属于任何已知的普适性类。CTMRG 结合有限 $\chi$ 标度分析可能给出答案。

\textbf{2. 更好的截断方案}：
是否存在利用 $N$ 皇后约束的\textbf{特殊结构}（如排列性质、对角线约束的单调性）的定制截断方案，使得有效键维度比朴素截断更小？

\textbf{3. MERA 表示}：
多尺度纠缠重整化（MERA）可以高效表示具有长程纠缠的态。$N$ 皇后问题的约束涉及多个尺度（短对角线到长对角线），MERA 的层次结构可能与之匹配。但如何具体构造 $N$ 皇后的 MERA 仍是开放问题。

\textbf{4. 量子计算的可能性}：
$N$ 皇后的张量网络结构可以映射到量子电路。量子计算机能否利用量子叠加和干涉来加速计数？Grover 搜索给出平方加速（$O(\sqrt{Z_N})$ 查询），但是否存在更好的量子算法？

\textbf{5. 推广到其他棋子}：
\begin{itemize}
\item 车（Rook）格点气体：只有行和列约束，没有对角线——这等价于计算矩阵的 permanent，张量网络是平面的
\item 象（Bishop）格点气体：只有对角线约束——同样可以构造张量网络，但拓扑不同
\item 国王（King）格点气体：约束是最近邻的——最简单的情况，等价于硬方块模型
\end{itemize}

这些变体模型提供了一个从简单到复杂的谱系，可以系统地研究约束结构如何影响张量网络方法的效率。

\subsection{更广阔的视角}

$N$ 皇后问题的张量网络分析为我们提供了一个独特的窗口，连接了三个通常被分开讨论的领域：

\begin{itemize}
\item \textbf{组合数学}：$N$ 皇后问题是经典的组合问题，精确计数属于 \#P-完全
\item \textbf{统计物理}：皇后格点气体是一个硬核粒子模型，具有丰富的相变行为
\item \textbf{量子信息}：张量网络的纠缠结构和计算复杂度有深刻的联系
\end{itemize}

张量网络方法的独特价值不在于替代传统的回溯搜索或蒙特卡洛方法，而在于提供了一种\textbf{统一的语言}和\textbf{系统化的近似框架}，使得组合问题可以用物理学家熟悉的工具来分析。

从实际应用的角度，最有前景的方向是用 CTMRG 研究皇后格点气体的相变行为——这是蒙特卡洛方法在高密度区域无法触及的领域，而张量网络方法可以给出可靠的定量结果。


%=============================================================
% 参考文献
%=============================================================
\begin{thebibliography}{30}

\bibitem{Orus2014} R.~Or\'{u}s, ``A practical introduction to tensor networks: Matrix product states and projected entangled pair states,'' \textit{Ann.~Phys.}~\textbf{349}, 117 (2014).

\bibitem{Schollwock2011} U.~Schollw\"{o}ck, ``The density-matrix renormalization group in the age of matrix product states,'' \textit{Ann.~Phys.}~\textbf{326}, 96 (2011).

\bibitem{LevinNave2007} M.~Levin and C.~P.~Nave, ``Tensor renormalization group approach to two-dimensional classical lattice models,'' \textit{Phys.~Rev.~Lett.}~\textbf{99}, 120601 (2007).

\bibitem{Xie2012} Z.~Y.~Xie \textit{et al.}, ``Coarse-graining renormalization by higher-order singular value decomposition,'' \textit{Phys.~Rev.~B}~\textbf{86}, 045139 (2012).

\bibitem{NishinoOkunishi1996} T.~Nishino and K.~Okunishi, ``Corner transfer matrix renormalization group method,'' \textit{J.~Phys.~Soc.~Jpn.}~\textbf{65}, 891 (1996).

\bibitem{Baxter1968} R.~J.~Baxter, ``Dimers on a rectangular lattice,'' \textit{J.~Math.~Phys.}~\textbf{9}, 650 (1968).

\bibitem{Baxter1982} R.~J.~Baxter, \textit{Exactly Solved Models in Statistical Mechanics} (Academic Press, 1982).

\bibitem{Kourtis2019} S.~Kourtis, C.~Chamon, E.~R.~Mucciolo, and A.~E.~Ruckenstein, ``Fast counting with tensor networks,'' \textit{SciPost Phys.}~\textbf{7}, 060 (2019).

\bibitem{GarciaSierra2021} D.~Garc\'{i}a-Mart\'{i}n and G.~Sierra, ``Quantum computation of the classically intractable permanent using tensor networks,'' \textit{Phys.~Rev.~A}~\textbf{104}, 032417 (2021).

\bibitem{Murg2005} V.~Murg, F.~Verstraete, and J.~I.~Cirac, ``Efficient evaluation of partition functions of frustrated and inhomogeneous spin systems,'' \textit{Phys.~Rev.~Lett.}~\textbf{95}, 057206 (2005).

\bibitem{Simkin2021} M.~Simkin, ``The number of $n$-queens configurations,'' (2021), arXiv:2107.13460.

\bibitem{Toda1991} S.~Toda, ``PP is as hard as the polynomial-time hierarchy,'' \textit{SIAM J.~Comput.}~\textbf{20}, 865 (1991).

\bibitem{OEIS} J.~Zhao, ``Queens on board,'' OEIS A000170.

\bibitem{Evenbly2015} G.~Evenbly and G.~Vidal, ``Tensor network renormalization,'' \textit{Phys.~Rev.~Lett.}~\textbf{115}, 180405 (2015).

\bibitem{Hauru2018} M.~Hauru, C.~Delcamp, and S.~Mizera, ``Renormalization of tensor networks using graph-independent local truncations,'' \textit{Phys.~Rev.~B}~\textbf{97}, 045111 (2018).

\end{thebibliography}

\end{document}
